% Acurácia do consenso ponderado por rodada, com conluio, 20 nós, EMA.
% Uma curva por percentual de maliciosos.
% O cenário de 0% não tem variante de conluio: não há nós maliciosos.
% Ajuste \figw e \figh abaixo.

\renewcommand{\figw}{0.92\textwidth}
\renewcommand{\figh}{7.6cm}

\begin{figure}[htbp]
\centering
\begin{tikzpicture}
\begin{axis}[tcc linha, width=\figw, height=\figh, ylabel={Acurácia}]
  \linhasDeTeste
  \addplot[serie zero] table[x=rodada, y=acuracia_ponderada, col sep=comma]
    {por_rodada/n20_m00_sem_ema.csv};
  \addlegendentry{0\%}
  \addplot[serie vinte] table[x=rodada, y=acuracia_ponderada, col sep=comma]
    {por_rodada/n20_m25_com_ema.csv};
  \addlegendentry{25\%}
  \addplot[serie cinquenta] table[x=rodada, y=acuracia_ponderada, col sep=comma]
    {por_rodada/n20_m50_com_ema.csv};
  \addlegendentry{50\%}
  \addplot[serie sessenta] table[x=rodada, y=acuracia_ponderada, col sep=comma]
    {por_rodada/n20_m65_com_ema.csv};
  \addlegendentry{65\%}
  \legendaRodadaDeTeste
\end{axis}
\end{tikzpicture}
\caption{Acurácia do consenso ponderado por rodada, com conluio, 20 nós, EMA.
Média entre as sementes. A curva de 0\% é o cenário sem maliciosos.}
\label{fig:acuracia-conluio}
\end{figure}
